\documentclass{article}
\usepackage[margin=2cm]{geometry}
\usepackage{amsmath}
\usepackage{graphicx}
\usepackage{booktabs}
\usepackage[utf8]{inputenc}
\usepackage{ragged2e}
\setlength{\emergencystretch}{3em}
\begin{document}
\sloppy

\subsection*{0.0.1 Summary of first-order equations}


To summarize, the first-order equations for the warp factor f and the scalars o,\$,\$3 are.
found to be


\begin{equation}
\begin{array} { c c } { f ^ { \prime } = 2 \left ( G _ { 0 } S _ { 0 } + C _ { 3 } S _ { 3 } \right ) } \\ { \sigma ^ { \prime } = 2 \eta \sqrt { N _ { 0 } ^ { 2 } + N _ { 3 } ^ { 2 } } } \\ { \cos \phi ^ { 3 } \left ( \phi ^ { 0 } \right ) ^ { \prime } = - \left ( G _ { 0 } M _ { 0 } + G _ { 5 } M _ { 3 } \right ) } \\ { = i \left ( G _ { 3 } M _ { 0 } + G _ { 5 } M _ { 3 } \right}\end{array}
\end{equation}


Furthermore, for consistency these were required to satisfy the algebraic constraint.


\begin{equation}
\begin{array} { c c } { e ^ { - 2 f } = 4 \left ( G _ { 0 } S _ { 0 } + G _ { 3 } S _ { 4 } \right ) ^ { 2 } - 4 \left ( S _ { 0 } ^ { 2 } + S _ { 3 } ^ { 2 } \right ) } & { } \\ {. } & {. } & { ( 2 ) } \end{array}
\end{equation}


The various functions featured in these equations were defined in and ..


\subsection*{0.1Numeric solutions}


In order to get acceptable numerical solutions from these equations, we must choose appropriate initial conditions. It is easy to check that the following initial conditions ensure
smoothness of all three scalars, as well as the vanishing of e2f at the origin,.


\begin{equation}
\begin{array} { c c } { \phi _ { 0 } ^ { 3 } = \sin ^ { - 1 } \left [ \frac { 1 } { 8 \operatorname { t a n h } \phi _ { 0 } ^ { 6 } } \left ( - 3 + \sqrt { 9 + 1 6 \operatorname { t a n h } ^ { 2 } \phi _ { 0 } ^ { 0 } } \right ) \right ] } & { } \\ { } & { } \\ { \sigma _ { 0 } = \frac { 1 } { 4 } \log \left [ \frac { \cosh \phi _ { 0 } ^ { 6 } \left ( 5 + \sqrt { 9 + 1 6 \operatorname { t a n h } ^ { 2 } \phi _ { 0 }}}}\end{array}\right.\right.
\end{equation}


We have defined for notational convenience \$8 = \$(0) and oo = o(0). For these initial.
conditions to be real, we must ensure that


\begin{equation}
\begin{array} { c } { | f ( \phi _ { 0 } ^ { 9 } ) | \leq 1 \qquad \qquad f ( \phi _ { 0 } ^ { 0 } ) \equiv \frac { 1 } { 8 \operatorname { t a n h } \phi _ { 0 } ^ { 6 } } \left ( - 3 + \sqrt { 9 + 1 6 \operatorname { t a n h } ^ { 2 } \phi _ { 0 } ^ { 6 } } \right ) \qquad \qquad ( 4 ) } \\ { \mathrm { o t } } \end{array}
\end{equation}


Noting that


\begin{equation}
\operatorname* { l i m } _ { \stackrel { \phi _ { 0 } ^ { 0 } \rightarrow - \infin } { * } } f ( \phi _ { 0 } ^ { 0 } ) = - \frac { 1 } { 4 } \qquad \qquad \quad \underset { a _ { 0 } ^ { 0 } \rightarrow + \infin } { \operatorname* { l i m } } f ( \phi _ { 0 } ^ { 0 } ) = \frac { 1 } { 4 } \qquad ( 5 )
\end{equation}


and also that f(8) is monotonically increasing, i.e.


\begin{equation}
\frac { d f } { d \phi _ { 0 } ^ { 0 } } > 0 ~
\end{equation}

\end{document}
